\section{Transformer 怎样真正学会预测下一个 Token？}
\label{intro:22c-language-training}

阅读指南提出了一种可大量练习的办法：\textbf{用原始文字本身作为后续 Token 的答案}。
前面已经学会网络如何计算分数，也学会 \conceptlink{loss}{损失函数}、\conceptlink{backprop}{反向传播} 和\conceptlink{gradient-descent}{梯度下降}。
现在才具备足够工具回答：\textbf{一整段真实训练文本，怎样变成位置对齐的预测目标、损失和参数更新？}
这里讲的是\textbf{一种自回归预训练任务}，不是说所有语言模型都只能这样学。

\subsection{“错开一格”：一段文本同时产生许多训练目标}
\phantomsection\label{concept:shifted-targets}

令完整 Token 序列为
\[
(t_1,t_2,t_3,t_4)=(1,3,4,2).
\]
把前面三项作为输入 $X_{\rm id}=(1,3,4)$，
把后面三项作为目标 $Y_{\rm id}=(3,4,2)$：
\[
\begin{array}{c|ccc}
\text{输入位置 }i&1&2&3\\\hline
\text{当前 Token}&1&3&4\\
\text{正确的下一 Token}&3&4&2
\end{array}
\]
这是\textbf{自监督训练}的一种典型形式：标签来自文本本身，
不需要额外人工为每个位置写一个分类标注。

\subsection{为什么训练时不会偷看未来？}

假设 $X$ 有三个 Token、形状 $3\times d_{\rm model}$。
如果允许第一个位置通过 Attention 读到第三个位置，
就可能让它在预测“第二个 Token”时偷看未来。
所以解码器使用\conceptlink{causal-mask}{因果掩码}：
\[
C=\begin{bmatrix}
0&-\infty&-\infty\\
0&0&-\infty\\
0&0&0
\end{bmatrix}.
\]
把它加到注意力分数上，再\textbf{逐行 Softmax}：
\begin{equation}
A=\softmax\!\left(\frac{QK^\top}{\sqrt{d_k}}+C\right).
\label{eq:masked-language-model}
\end{equation}
第一个 Query 只能看见第一个 Key；第二个可看见前两个；
第三个可以看见前三个。图中的 $-\infty$ 是数学上的理想化表达，
实现中会使用正确处理遮蔽位置的数值方案。

\begin{idea}[title={为什么“训练并行”和“生成逐词”并不矛盾？}]
训练时，整段文本的正确输入都已经给出。
掩码确保每个位置只使用允许看到的历史信息，
所以多个位置的预测可以在一次矩阵计算中并行得到。

生成时，下一个真实 Token 还不知道，必须先产生一个，
把它追加到序列末尾，才能继续生成下一个。
这叫\textbf{自回归生成}。因此\textbf{训练的并行性}
不等于\textbf{解码时所有未来 Token 可以同时确定}。
\end{idea}

\subsection{输出张量和目标怎样对齐？}
\phantomsection\label{concept:cross-entropy}

设词表大小为 $|\mathcal V|=5$，
模型输出\conceptlink{logits}{logits}：
\[
Z\in\R^{3\times5}.
\]
三行分别给出三个位置对词表中五个候选的分数。
逐行经过 Softmax，得到预测概率 $P_{i,j}$。

\textbf{先把上一章的结果接进来：}
上一章用“春天、花”预测下一 Token。
真实答案是 ID 4“开了”，但初始模型只给它约 $0.059$ 的概率。
这道题的负对数损失因此是
\[
L_{\text{这一题}}=-\log(0.05923)\approx2.826.
\]
如果把正确候选的概率提高，损失就会降低；
因此\textbf{从候选概率到训练信号}，只差一步可以求导的损失。

第 $i$ 个位置的交叉熵损失取\textbf{正确 Token 所在列}：
\begin{equation}
L_i=-\log P_{i,t_{i+1}}.
\end{equation}
整个序列的平均损失是：
\begin{equation}
\boxed{L=-\frac{1}{3}\sum_{i=1}^3\log P_{i,t_{i+1}}.}
\end{equation}
接下来\textbf{换一组人为设定的三位置概率}，演示怎样合并整段文本的损失；
它不再使用上一章那组初始输出概率。
如果三次正确 Token 的预测概率分别是
$0.5$、$0.25$、$0.5$，则
\[
L=-\frac{\log0.5+\log0.25+\log0.5}{3}
=\frac{4\log2}{3}\approx0.924.
\]
注意这是\textbf{一个人工算例}，不是实际模型的测试成绩。

\subsection{亲手训练：让“开了”超过“花”}
\label{sec:transformer-gradient-step}

沿用上一章输出的隐藏行向量
$\vect h_{\rm last}^\top=(0.495,2.505)$，
初始 logits 是 $(0,0.495,0,2.505,0)$。
正确 ID 4 的概率只有约 $0.059$，而错误 ID 3 为 $0.725$。
令真实 One-hot 标签为 $\vect y_{\rm target}^\top=(0,0,0,0,1)$，
交叉熵的 logit 梯度为
\[
\boxed{\frac{\partial L}{\partial\vect z^\top}
=\vect p^\top-\vect y_{\rm target}^\top
\approx(0.059,0.097,0.059,0.725,-0.941).}
\]
可见错误类别 ID 3 要往下调，正确类别 ID 4 要往上调。
因为 $\vect z^\top=\vect h_{\rm last}^\top W_{\rm vocab}$，
词表权重梯度是外积：
\begin{equation}
\boxed{\frac{\partial L}{\partial W_{\rm vocab}}
=\vect h_{\rm last}(\vect p-\vect y_{\rm target})^\top
\in\R^{2\times5}.}
\end{equation}

\begin{example}
取学习率 $\eta=0.1$，同时更新输出头五列。第一步
$W\leftarrow W-\eta\partial L/\partial W$，
其中错误的 ID 3 列与正确的 ID 4 列分别变为
\[
\begin{aligned}
W_{:,3}&:(0,1)^\top\ \mapsto\ (-0.0359,0.8183)^\top,\\
W_{:,4}&:(0,0)^\top\ \mapsto\ (0.0466,0.2357)^\top.
\end{aligned}
\]
每次重新计算概率和梯度，得到以下数值：
\begin{center}\small
\begin{tabular}{@{}ccccc@{}}\toprule
更新次数 & $p(\text{花})$ & $p(\text{开了})$ & 损失 $-\log p_4$ & 预测\\\midrule
0 & 0.725 & 0.059 & 2.826 & 花（错）\\
1 & 0.590 & 0.143 & 1.947 & 花（错）\\
2 & 0.445 & 0.276 & 1.286 & 花（错）\\
3 & 0.321 & 0.427 & 0.851 & 开了（对）\\\bottomrule
\end{tabular}
\end{center}
\end{example}
三次更新后概率排序翻转。这是\textbf{固定隐藏状态，只训练输出头}的局部实验，
不是端到端训练已全部完成的证据。

\subsection{再走一步：梯度怎样进入 FFN 和 Embedding？}
现在不再把输出头当作反向传播的终点。
用\textbf{初始}的输出头权重，隐藏状态的梯度是
\[
\underbrace{\frac{\partial L}{\partial\vect h_{\rm last}^\top}}_{1\times2}
=(\vect p-\vect y_{\rm target})^\top W_{\rm vocab}^\top
\approx\boxed{(0.097,0.725)}.
\]
本例 FFN 为 $0.5\relu$，最后位置两个特征都为正，
于是 $X_{{\rm block},2,:}=1.5X_{{\rm att},2,:}$，
反传到 FFN 输入（含残差）的梯度是
\[
\frac{\partial L}{\partial X_{{\rm att},2,:}}
=1.5(0.097,0.725)\approx\boxed{(0.146,1.088)}.
\]
这份梯度既能沿残差恒等支路回到输入 Embedding，
也会经过 Attention 的 Value 读出与 Query/Key Softmax 的分数路径。
多条路径对同一参数的梯度需要\textbf{求和}，
才能更新 FFN、Q/K/V 和相关 Embedding 行。
\subsection{真正向前层更新：同时改 FFN 和 Value 投影}

前面三轮 SGD 只训练了输出头。现在保留\textbf{初始}的输出头，
沿同一计算图再走一个\textbf{不同的单步实验}：
同时更新 FFN 的下投影 $W_{\rm down}$ 与 Attention 的 $W_V$。
所有梯度都从\textbf{更新前}的权重计算，不混用新旧参数。

第二个位置在残差后为
$u=(0.330,1.670)$，ReLU 后仍是 $u$；
此时 $W_{\rm down}=0.5I_2$，
$h=u+\relu(u)W_{\rm down}=(0.495,2.505)$。
从词表头反传的梯度
$g_h=\partial L/\partial h=(0.09716,0.72516)$。
因此下投影参数梯度为
\[
\frac{\partial L}{\partial W_{\rm down}}
=\underbrace{u^\top}_{2\times1}
\underbrace{g_h}_{1\times2}
\approx
\begin{bmatrix}
0.03206&0.23930\\
0.16226&1.21101
\end{bmatrix}.
\]
为了继续传回 Attention，FFN 加残差对 $u$ 的导数
在本例两个正分量处是 $I_2+W_{\rm down}=1.5I_2$，
所以 $g_u=1.5g_h\approx(0.14574,1.08773)$。
上一章给出的注意力行是 $a=(0.330,0.670)$，
原输入矩阵 $X=I_2$；保持 Q/K 不变，
$O_{\rm att,last}=aXW_V=aW_V$。
链式法则再次给出一个\textbf{真正可训练参数}的梯度：
\[
\frac{\partial L}{\partial W_V}
=\underbrace{a^\top}_{2\times1}
\underbrace{g_u}_{1\times2}
\approx
\begin{bmatrix}
0.04810&0.35895\\
0.09765&0.72878
\end{bmatrix}.
\]
让两组矩阵\textbf{同时}按学习率 $0.1$ 更新，
重新前向计算最后位置（所有其它矩阵不变），得到
\[
W_V'\approx\begin{bmatrix}
0.99519&-0.03590\\-0.00976&0.92712
\end{bmatrix},\quad
W_{\rm down}'\approx\begin{bmatrix}
0.49679&-0.02393\\-0.01623&0.37890
\end{bmatrix}.
\]
\[
h_{\rm last}'\approx(0.45566,2.21140),\qquad
p_{\rm true}':0.05923\longrightarrow\boxed{0.07296},
\qquad L:2.826\longrightarrow\boxed{2.618}.
\]
这仍然只是单样本单步训练；但梯度确实经过 FFN、残差、
Attention 的 Value 分支，修改了\textbf{两组中间层权重}并降低损失。
Q/K 和 Embedding 的梯度仍可依同样规则继续推导，
不能把本例称作训练了全部参数。

\subsection{哪些权重会因此而更新？}

这份损失从词表输出头往回，经过每一层 Transformer Block，
一路传到 Q/K/V 投影、FFN、归一化参数和 Embedding 表：
\[
\text{损失}\ \longrightarrow\
\text{反向传播}\ \longrightarrow\
\nabla_{\theta}L\
\longrightarrow\
\text{优化器更新 }\theta.
\]
$\theta$ 表示全部参与训练的可学习参数；
实际是否更新某组参数，取决于它是否被冻结、是否参与计算以及训练设置。
同一参数在多个 Token 位置上被复用，
它的梯度要汇总这些路径的贡献。

\subsection{训练时为什么有时需要忽略一些 Token？}

真实批次可能含有填充（padding）或不需要参与训练的提示片段。
因此实现中常定义\textbf{损失掩码}：仅对选定的目标位置累计损失。
它与\textbf{因果注意力掩码}不是一回事：
前者决定哪里\textbf{计分}，后者决定哪里\textbf{可读取}。

\begin{warningbox}
从下一个 Token 的损失最小化，
不能直接推出模型会在所有任务上做对事实判断，
也不能推出它已经学会遵循人类指令。
后面的监督微调、偏好优化与系统工具能力解决的是不同层次的问题。
\end{warningbox}

\begin{conclusion}
对于本书的自回归训练任务，\textbf{错位标签}提供答案，\textbf{因果掩码}限制可见信息，
\textbf{逐位置概率和交叉熵}定义误差，\textbf{反向传播}推动参数学习。
但训练结束后，用户不会提供未来文本作为标签：
下一章研究\textbf{没有标准答案在旁边时，模型怎样一步步生成自己的回答}。
\end{conclusion}
